\subsection{Историко-философские истоки формирования политических ценностей: от Античности до Нового времени}

Генезис политической культуры как системы ценностных ориентаций, нормативных стандартов и образцов взаимодействия между властью и обществом имеет глубокие историко-философские корни. Формирование первых устойчивых представлений о месте индивида в политической системе началось в эпоху Античности, когда в фокус философской мысли впервые попали проблемы государственного устройства, гражданского долга и легитимности власти. 

В трудах Платона и Аристотеля были заложены основы понимания того, что устойчивость государства напрямую зависит от ценностных ориентиров и моральных качеств его граждан. Платон в диалоге <<Государство>> связал типы политических режимов с доминирующими чертами характера людей, их населяющих, фактически предвосхитив современные концепции субкультур. Фундаментальный вклад в развитие античных представлений о политической культуре внес Аристотель. В своем трактате <<Политика>> он определил человека как <<животное политическое>> (zoon politikon), подчеркнув, что полноценное существование индивида возможно только в рамках полиса\footnote{Аристотель. Политика // Сочинения: В 4 т. Т. 4. --- М.: Мысль, 1983. --- С. 379.}. Аристотель классифицировал формы правления на основе их ценностного содержания (служение общему благу или личной выгоде правящих) и выделил <<политию>> как наиболее стабильный режим, опирающийся на ценности умеренности и многочисленный средний класс. Таким образом, в античности сформировалась нормативно-активистская модель политических ценностей, где участие в делах общины признавалось высшим благом и обязанностью свободного гражданина.

В эпоху Средневековья характер политических ценностей претерпел радикальную трансформацию под влиянием религиозного догматизма и теоцентризма (труды Аврелия Августина, Фомы Аквинского). Политическая культура этого периода приобрела выраженный подданнический характер: земная власть рассматривалась как божественное установление, а ключевыми добродетелями подданных провозглашались смирение, послушание и конформизм.

Принципиальный поворот к секуляризации и реалистическому осмыслению политического поведения произошел в эпоху Возрождения благодаря работам Н. Макиавелли. Он полностью отделил политическую науку от морали и религии, начав исследовать политику такой, какая она есть в реальности, а не какой она должна быть. Н. Макиавелли описал специфику прагматической политической культуры правящей элиты, основанной на рациональном расчете, манипулировании и жестком удержании власти, что заложило основы для будущих исследований психологии электорального и элитарного поведения.

Новый этап эволюции политических ценностей связан с эпохой Нового времени и концепциями общественного договора (Т. Гоббс, Дж. Локк, Ж.-Ж. Руссо). В этот период происходит переход от концепции подданничества к идее гражданственности. В трудах Дж. Локка оформляются ценности классического либерализма: неотчуждаемость естественных прав человека (на жизнь, свободу и собственность), верховенство закона и ограничение государственного вмешательства. Ж.-Ж. Руссо, развивая идею народного суверенитета, актуализировал античный идеал активного соучастия граждан в управлении. Философия Нового времени сформировала ценностную матрицу современной гражданской политической культуры, провозгласив приоритет прав личности над государством и легитимность власти только при условии согласия управляемых.